% 本文件由 LaTeX 排版 Agent 根据有效 translation.json 自动生成。
% author=Tertullian; work_id=0320; title=论悔改

\chapter{论悔改}

% structure: p0001 source=h1 target=omitted reason=page-title-already-rendered-by-parent

% structure: p0002 source=h2 target=section reason=relative-heading-rank; base=h2

\section{第一章 论外邦人的悔改}

人们按本性所能理解，悔改是一种因厌恶某种\Emph{先前珍视的}较恶劣情感而产生的心灵悸动——我所指的是那种人，甚至我们自己在过去也一度如此：盲目，没有主的光明。然而，他们距悔改的\Emph{理性}之远，正如他们距理性本身之源之远。事实上，\Emph{理性}是天主的事，因为创造万有的天主没有不\Emph{借由理性}提供、布置、安排任何事物——祂也愿意一切事物借由\Emph{理性}来处理和理解。因此，凡不认识天主的，必然也不认识属于祂的事物，因为没有宝库是对陌生人开放的。如此，他们在没有理性之舵的情况下航行生命的全部路程，不知道如何躲避那即将降临于世界的飓风。此外，他们在实践悔改时表现得多么不合理，只需这一点就足以简要表明：他们甚至在\Emph{善}行上也实行悔改。他们悔改于信实、爱、纯朴、忍耐、仁慈，恰与任何\Emph{受这些情感驱使的}行为落在不知感恩之地成正比。他们因行了善而咒诅自己；并且将那种主要适用于最善之行的悔改铭记在心，用心记住绝不再行善。相反，对于\Emph{恶}行的悔改，他们却看得较轻。总之，他们把这同一（德行）当作更容易\Emph{犯罪}的手段，而非\Emph{行义}的手段。\par

% structure: p0004 source=h2 target=section reason=relative-heading-rank; base=h2

\section{第二章 真悔改是神圣之事，源自天主，并服从祂的律法}

然而，倘若他们行事如同那些有分于天主、因而也有分于理性的人，他们必会首先慎重权衡悔改的重要性，而绝不会用它作为\Emph{自我定罪}的理由，进行乖谬的自新。简言之，他们会限制悔改的尺度，因为他们也会达到犯罪的限度——我是指因敬畏天主。但是，哪里没有敬畏，同样也就没有改正；哪里没有改正，悔改就必是虚浮的，因为它缺少天主播种它所期望的果实，即人的救恩。因为天主——在人类轻率所犯的如此众多和如此重大的罪恶之后（始于人类之首亚当），在人被判罪、连同世界的陪嫁——在他被逐出乐园并屈服于死亡之后——祂急速返回自己的仁慈，从那时起就在祂自己内开创了悔改，撤销了祂初次震怒的判决，承诺宽恕祂自己的作品和肖像。于是，祂为自己聚集了一个民族，以祂恩惠的许多慷慨分配来培育他们，并且，在如此频繁地发现他们最忘恩负义之后，仍然不断劝勉他们悔改，并派遣全体先知的声音去预言。不久，祂自由地应许了在末世祂要借着圣神像光洪流一样倾注在全世界的恩宠，祂命令悔改的洗礼先行，目的是借着悔改的标记与封印，首先预备那些祂要通过恩宠召唤，去承受确切应许给\Person{亚巴郎}的（产业）的人。\Person{若翰}并不缄默，他说：\Quote{你们要进入悔改，因为现在救恩将临近万邦}——那就是主，照着天主的应许带来救恩。\Person{若翰}作为祂的先驱，把（他所宣讲的）悔改指向祂；这悔改的职责是洁净人的心灵，以便任何根深蒂固的错误所灌输的污秽，任何无知在人心中所滋生的污染，\Emph{那}悔改都应将其扫除、刮净、抛出门外，从而预备心的居所，使之洁净，迎接即将来临的圣神，使祂能愉快地带着祂天上的祝福进入其中。这些祝福的标题\Emph{简言之}是一个——人的救恩——以废除先前的罪为预备步骤。这是悔改的（最终）原因，这是她的工作，承接神圣慈悲的事业。凡有益于人的，就是服务天主。然而，我们认识主时所学习的悔改的\Emph{规则}保持一个确定的形式——\Emph{即}，绝不可对\Emph{善}行或善念施加强力（可说是暴力）。因为天主从不认可对\Emph{善}行的弃绝，因为它们是祂自己的（祂既是创造者，也必然是保护者），同样祂接纳它们，若是接纳者，也是赏报者。那么，让人的忘恩负义去注意吧，如果它把悔改也加诸于善工；让他们的感恩也去注意吧，如果赚取感恩的愿望是行善的动机：两者都是属世俗和必死的。因为你向感恩的人行善，收益何等小！向忘恩的人行善，损失何等大！一件\Emph{善}行以天主为债务人，同样一件\Emph{恶}行也有天主为债务人；因为审判者是每一案件的赏报者。既然如此，既然天主作为审判者主持着公义的执行与维护（这对祂最为珍贵），既然祂是以公义为着眼点来制定祂全部诫命的总和，那么，还有怀疑的余地吗？——正如我们在一切行为中普遍如此，在悔改的事上，也必须向天主呈献公义？这职责确实只能在一个条件下实现：悔改必须\Emph{只}应用在\Emph{罪}上。再者，没有行为除了\Emph{恶}行配称为\Emph{罪}，也没有人因行善而犯错。但若他没有犯错，为什么他要闯入悔改（的领域）——那是犯错者的私人领地？为什么他要将属于邪恶的职责强加于他的善良？结果就变成：当一件东西被用在不该用的地方，那么它该在的地方，就被忽视了。\par

% structure: p0006 source=h2 target=section reason=relative-heading-rank; base=h2

\section{第三章 罪可分为肉身之罪和属神之罪。两者同样受天主的审断与惩罚，即使不受人的审断。}

那么，究竟为哪些事悔改才算公正合宜——即应归入\Quote{罪}这一标题之下的有哪些——时机确实要求我加以记录；但这样做似乎不必要。因为当认识主时，我们的精神被自己的作者所\BibleQuote{回顾}\BibleRef{路 22:61}，便不由自主地进入对真理的认识；并得以认识天主的诫命，立即受其教导说：\Quote{凡天主命我们戒避的，应算为\Emph{罪}：}因为大家公认天主是某种至高的善的\Emph{本质}，当然只有恶才会使善不悦；因为，相互对立的事物之间，没有友谊。然而，简要地提及以下事实并非令人厌烦：即罪中，有些是肉欲的，即属于身体的；有些是精神的。因为人是由这两种实体的结合所组成，他罪恶的根源无非是他组成的根源。但并不是说身体和精神是两件事就构成罪的相互不同——否则它们反倒因此而更\Emph{相等}，因为\Emph{两者}组成\Emph{一}——免得有人按其实体的区别来区分它们的\Emph{罪}，以至认为一个比另一个更轻或更重：如果（确实如此）肉体和精神都是天主的受造物，一个由祂手所造，一个由祂的\Emph{气息}所完成。既然它们同样属于主，那么其中任何一个\Emph{犯罪}都同样\Emph{冒犯}主。难道你有权区分肉体和精神的行为吗？它们在生命、死亡和复活中的共融与结合如此密切，以致\Quote{在那时}它们被同样复活，或进入生命，或进入审判；因为，诚然，它们同样要么犯了罪，要么无罪生活？我们暂且一次性地预设这一点，以便我们明白：如果人任何一部分在任何事上犯了罪，则悔改加于那一部分的责任并不少于加于\Emph{两者}共同的责任。两者的\Emph{罪责}是共同的；\Emph{审判者}——即天主——也是共同的；因此，悔改的治疗之药也是共同的。罪被称为\Quote{精神的}和\Quote{肉身的}的根源在于：每一件罪要么是\Emph{行为}的，要么是\Emph{思想}的：所以，在\Emph{行为}中的罪是\Quote{肉身的}，因为\Emph{行为}像\Emph{身体}一样能被\Emph{看见}和\Emph{触摸}；在\Emph{心思}中的罪是\Quote{精神的}，因为\Emph{精神}既不被\Emph{看见}也不被\Emph{触摸}：由此表明，不仅行为的罪，而且意志的罪，都应当避免，并藉悔改洁净。因为人的有限性只审判行为的罪，由于它无法穿透意志的隐藏处，我们不要因此在天主眼中轻视意志的罪行。天主是全然充足的；任何罪的来源都逃不过祂的眼目；因为祂既不无知，也不忽略宣判审判。祂不掩饰也不欺骗自己敏锐的眼目。意志是行为的根源，这一点（我们该说什么）呢？因为若有些罪归因于偶然、必然或无知，就由他们自行负责；若排除这些，则除非出于\Emph{意志}，否则无犯罪。既然意志是行为的根源，那么它难道不因其在罪责中的首要地位而更应受惩罚吗？即使有些困难阻碍了它的完全实现，它也并不因此被免除责任；因为它自身归咎于自身：即使它在自己能力范围内完成了工作，也并不会因实现的失败而可原谅。事实上，主是如何表明祂为法律加建上层建筑，除非也禁止意志的罪（以及其他罪）？祂不仅定规实际侵犯他人婚姻的人是奸夫，同样也定规因贪欲的目光玷污（妇女）的人是奸夫。因此，心思将禁止做的事摆在自己面前，并鲁莽地通过意志完善其执行，是足够危险的。既然意志的能力如此之大，即使没有完全满足其私欲，也等同于行为；因此，它将作为行为受罚。说\Quote{我\Emph{愿意}但\Emph{没有}做}是完全虚妄的。相反，你\Emph{应当}完成那事，\Emph{因为}你愿意；否则就不愿意，因为你没有完成它。但是，藉着你良知的承认，你宣判了自己的定罪。因为如果你热切渴望一件\Emph{好}事，你就会急于完成它；同样，既然你没有完成一件\Emph{恶}事，你就不该热切渴望它。无论你持哪种立场，你都被罪责紧紧束缚；因为你要么\Emph{愿意}了恶，要么没有\Emph{完成}善。\par

% structure: p0008 source=h2 target=section reason=relative-heading-rank; base=h2

\section{第四章 悔改适用于各类罪。不仅要、也不主要为了其带来的益处而实行，而是因为天主的命令。}

因此，凡由肉身或精神、由行为或意志所犯的一切罪，那位藉审判定下了刑罚的同一\Emph{天主}，也藉悔改应许了宽恕，向百姓说：\Quote{你们悔改吧，我必拯救你们；}又说：\Quote{我指着我的永生起誓——主说——我宁愿（有）悔改，而不愿死亡。}悔改因此是\Quote{生命}，因为它被置于\Quote{死亡}之上。你这罪人同我一样（不，比我更轻，因为我在罪中的优越地位，我承认是我的），你要如此急切地奔向悔改，如此拥抱它，如同一个遭遇海难的人抱住某块木板。当你沉没在罪的波涛中时，这木板将把你拉出，并将你载入天主仁慈的港口。抓住这意想不到的幸福时机：你从前在天主眼中不过是\BibleQuote{桶里的一滴水}\BibleRef{依 40:15}和\Quote{禾场上的灰尘}以及\Quote{窑匠的器皿}，从今以后却要成为\Quote{那栽种在水边、叶子常青、按时结果}的树，且看不到\BibleQuote{火}或\BibleQuote{斧头}。\BibleRef{玛 3:10}既已找到\BibleQuote{真理}\BibleRef{若 14:6}，就当为错误悔改；为爱上天主所不爱之物而悔改：就连我们也不允许我们的奴仆去恨那些冒犯我们的事物；因为自愿服从的原则在于心思的相似。\par

要数算悔改的益处，题材极为丰富，因此应交给伟大的口才。然而，我们按自己狭隘的能力，只强调一点：天主所命定的，是善且最好的。我认为，争辩神圣诫命之\Quote{善}是僭越；因为事实上，并非因为它是善的才迫使我们去服从，而是因为天主已命定了它。要索取服从，神圣权能的威严拥有优先权利；命令者的权威先于服役者的益处。\Quote{悔改是善还是不善？}你为何踌躇？天主命令；不，祂不仅命令，还劝勉。祂藉赏报——即救恩——来邀请，甚至藉起誓说\Quote{我活着}，希望我们相信祂。哦，我们有福了，天主竟为我们起誓！哦，我们最悲惨了，若连主起誓我们都不信！因此，天主如此高度称赞、甚至（按人的方式）以誓言作证的事，我们当然必须趋近，并以最大的郑重守护；好使我们恒久地持守在神圣恩宠的庄严保证的（信德）中，从而也能同样地持守在它的果实和益处中。\par

% structure: p0011 source=h2 target=section reason=relative-heading-rank; base=h2

\section{第五章 悔改后决不可再重返于罪}

因为我要说的是：那藉天主的恩宠显示并命令给我们的悔改，它将我们召回与主同在的恩宠中，一旦我们学会并采用了它，此后决不应因重复犯罪而取消。现在，再也无可推诿的借口为你辩护了；你在认识了主并接受了祂的诫命之后——简言之，在着手为（过去的）罪而悔改之后——竟然又再投身于罪。这样，你离无知越远，就与顽抗结合得越紧。因为如果你曾悔改犯罪，乃是因为你开始敬畏主，那么为何你竟宁愿取消你因敬畏而做的事？除非你已不再敬畏。因为除了顽抗，没有任何事能颠覆敬畏。既然连那些不认识主的人，也因对主的无知——因为天主公开地摆在人面前，甚至从祂天上的恩惠来看也是可以理解的——而得不到免除刑罚的例外，那么已知悉祂时却轻视祂，这是多么危险啊！而轻视祂的人，正是那些藉祂的帮助认识了善恶之事，却常常冒犯自己的理智——即天主的恩赐——通过重拾他所理解应当避免且已避免过的事；他抛弃了恩赐，也就拒绝了赐予者；他否定了恩主，因不尊崇恩惠。他怎能取悦祂，若那恩赐自己都不悦？这样，他不仅表现出对主的顽抗，更显为忘恩负义。此外，那人向主所犯的罪不轻，因为他曾藉悔改弃绝了主的仇敌魔鬼，并在此名下将之屈服于主，却又藉自己的回归（给仇敌）再度高举他，使自己成为他的欣喜；以至那恶者因得回猎物而重新向主夸耀。他岂不是——这样说甚至危险，但为了益处必须提出——将魔鬼置于主前？因为他似乎作了比较，认识了二者；并且公正地判决那一位更好，他宁愿再做那一位的（仆人）。这样，那曾藉罪之悔改开始向主赔补的人，又将藉对自己悔改的第二次悔改向魔鬼赔补，并且他越讨主的仇敌喜悦，就越被天主憎恨。但有人说：\Quote{天主若以心灵和思想仰望祂，即使未以\Emph{外在}行为实行，祂也满意；这样他们犯罪却无损于他们的敬畏和信德。}这就是说，他们犯奸淫却无损于他们的贞洁；他们为父母下毒却无损于他们的孝道！那么，他们自己也将被投入地狱却无损于他们的赦免，因他们犯罪却无损于他们的敬畏！这是悖逆的首要例证：他们犯罪，是因为他们敬畏！我猜想，如果他们不敬畏，就不会犯罪！因此，谁若不愿天主被冒犯，就根本不要敬畏祂，如果敬畏是犯罪的理由。但这些性情惯常从伪善者的种子生发出来，他们与魔鬼的友谊是不可分割的，他们的悔改从不忠实。\par

% structure: p0013 source=h2 target=section reason=relative-heading-rank; base=h2

\section{第六章 切勿擅自领受圣洗圣事。须有先行的悔改，并由生活的改正来显明}

那么，我们贫弱的能力试图就一次彻底把握悔改并永久保持它而提出的一切，确实关系到所有归向于主的人，因为他们都是为获得天主的恩宠而竞争救恩的人；但尤其迫切的是那些年轻的初学者，他们刚刚开始用天主的言论浸润耳朵\BibleRef{申 32:2}，如尚在早期幼年的幼犬，眼睛尚未完全睁开，摇摇晃晃地爬行，嘴里说放弃以前的行为，并接纳（悔改的）宣认，却忽略了完成它。因为欲望的最终目标驱使他们渴望一些以前的\Emph{行为}，正如水果，当它们已经开始变得酸涩或苦老时，仍在某些部分取悦自己的可爱。此外，对圣洗圣事的自以为是的确信，引入了各种恶性的拖延和推诿，关于悔改；因为，确信自己的罪必得赦免，\Emph{人们}同时窃取了中间的时间，将其变成自己犯罪的美好假期，而不是学习不犯罪的时间。再者，希望赦罪赐予一个他们尚未完成的悔改，这是多么不一致！这是伸手要商品，却不交出代价。因为悔改是主决定赐予赦免的代价：祂提出以悔改的补偿交换来赎买脱离惩罚。那么，如果卖主先检查他们交易用的钱币，看是否被切割、刮擦或掺杂，我们同样相信，主在即将赐予我们如此昂贵的商品，即永生时，先对我们的悔改进行考验。\Quote{但与此同时，让我们推迟我们悔改的实际；那么，我想，当我们被赦免时，就会清楚我们已经改正了。}绝不；（但我们的改正应当显明）当赦免悬而未决，仍面临惩罚时；当\Emph{悔改的人}尚未配得——就我们所能配得而言——他的释放时；当天主在威胁时，而不是在宽恕时。因为哪个奴隶，在获得自由改变地位后，会责备自己的（过去）偷窃和逃离？哪个士兵，在退役后，会为自己的（以前）烙印进行弥补？罪人必须在接受赦免之前哀哭自己，因为悔改的时间与危险和恐惧的时间一致。我并非否认，天主的恩惠——我指的是罪的除去——对即将进入（圣洗的）水的人肯定有效；但我们必须努力的是，使我们得蒙恩赐获得那福佑。因为谁会赐予你，一个如此不忠信悔改的人，任何水的一次洒水？的确，悄悄接近它，并让主管此事的管理员被你的断言误导，是容易的；但天主为祂自己的宝藏预作准备，不容许不配的人偷袭它。实际上，祂说什么？\BibleQuote{没有隐藏的事，不會不被显露出来的。}\BibleRef{路 8:17}随你用什么黑暗（帷幕）掩盖你的行为，\BibleQuote{天主是光。}\BibleRef{若一 1:5}但有些人以为天主仿佛有义务将祂承诺（要赐予）的甚至赐给不配的人；他们就把祂的慷慨变成了奴役。但如果天主出于必要赐给我们死亡的标记，那么祂这样做是出于\Emph{不情愿地}。但谁允许一个不情愿赐予的礼物被永久保留呢？因为许多人岂非后来失落了（恩宠）？这礼物难道不是从许多人那里被夺走吗？这些，无疑，就是那些偷袭（宝藏）的人，他们在接近悔改的信德之后，在沙土上建造注定倒塌的房子。那么，没有人应该因被归入学习者的\Quote{新兵班}而自夸，仿佛因此他现在就有犯罪的许可。一旦你\Quote{认识主}，你就应当敬畏祂；一旦你凝视了祂，你就应当尊敬祂。但\Emph{你的}\Quote{认识}祂又有什么不同，当你停留在与过去的日子相同的行径中，那时你\Emph{不}认识祂？再者，是什么将你与一位成全的天主仆人区分开来？难道有一个基督是为已受洗者，另一个是为学习者？难道他们有不同希望或赏报？不同的审判恐惧？不同的悔改必要性？那\Emph{圣洗的}洗涤是信德的封印，这信德始于悔改的信德并受其推荐。我们受洗不是\Emph{为了}我们\Emph{可以}停止犯罪，而是\Emph{因为}我们\Emph{已经}停止，因为在\Emph{心里}我们\Emph{已经}\Emph{被}洗净了。因为学习者的\Emph{第一次}圣洗就是\Emph{这}：完美的敬畏；此后，就你对主的理解而言，信德\Emph{是}健全的，良心一次而永远地拥抱了悔改。否则，如果只是经过圣洗的水之后我们才停止犯罪，那么我们是出于\Emph{必要}，而非出于\Emph{自由意志}，才穿上无辜。那么，谁在善上更卓越？是那不被\Emph{允许}作恶的人，还是那\Emph{厌恶}作恶的人？是那被\Emph{命令}无罪的人，还是那\Emph{以之为乐}的人？那么，让我们不要停止偷窃的手，除非门闩的坚硬阻止我们；也不约束眼睛不去看淫乱的贪欲，除非有监护人的管束；如果没有任何一个将自己交付于主的人，除非被圣洗约束，会停止犯罪。但如果有人持这种看法，我不知道他受洗后，是否会因想到自己已经\Emph{停止}犯罪而感到更多悲伤，而不是因已经\Emph{逃脱}而喜乐。因此，学习者应当\Emph{渴望}圣洗，但不仓促\Emph{接受}它：因为渴望它的人，尊敬它；仓促接受它的人，藐视它：前者显出谦逊，后者显出傲慢；前者满足它，后者忽略它；前者渴望赢取它，后者却把它当作应得的回报许诺给自己；前者接受，后者僭取。你会判断谁更值得，除非是更改正的人？谁是更改正，除非是更畏惧的人，并因此履行了真正悔改的责任？因为他害怕仍然留在罪中，以免他不配得\Emph{圣洗的}接纳。但仓促接受者，既然他许诺了自己（作为应得），自然确信（自己会得到），就\Emph{不能}害怕：因此他也没有履行悔改，因为他缺乏悔改的工具因，即恐惧。仓促接受是不敬虔的份；它使寻求者膨胀，它藐视施与者。因此它有时欺骗，因为它在礼物到期之前许诺自己\Emph{礼物}；从而那位要提供\Emph{礼物}的主常被冒犯。\par

% structure: p0015 source=h2 target=section reason=relative-heading-rank; base=h2

\section{第7章 论已领受圣洗后失足者的补赎}

主基督，但愿祢的仆人蒙受恩赐，得以学习或聆听关于补赎之训诲，正如他们身为\Emph{学习者}时所当行的，即不再犯罪；换言之，但愿他们此后一无所知于补赎，\Emph{并且}无需乎此。提及\Emph{第二种}——实乃\Emph{最后}——希望，实令人不快；以免因论述尚存之补救性悔改，而似乎指向更多犯罪的空间。但愿无人如此曲解我们的意思，\Emph{仿佛}因有补赎之门，便因此亦有犯罪之门；又仿佛天上仁慈的丰盛成了人类妄为的许可。但愿无人因天主更仁慈而变得不那么良善，屡次犯罪如同屡次被赦免。否则，他必将发现\Emph{逃避}的终结——当他找不到\Emph{犯罪}的终结时。我们已逃脱了\Emph{一次}：至此为止，\Emph{不可再越雷池}，即使我们似乎可能再次逃脱。人在遭遇海难后，往往从此与船和海洋决绝；并因\Emph{珍视}对危险的记忆，而尊崇天主所赐的\Emph{恩惠}——即他们的得救。我赞美他们的畏惧，我喜爱他们的敬畏；他们不愿再次成为天主仁慈的负担；他们害怕似乎践踏已获得的恩惠；他们以至少是良好的谨慎，避免再次尝试那曾使其畏惧之事。因此，他们妄为的限度便是他们畏惧的明证。\par

况且，人的畏惧乃是天主的荣耀。但无论如何，那（我们）最顽固的仇敌从不使他的恶意有喘息之机；事实上，当他充分感到一个人已脱离\Emph{他的魔掌}时，他最为残忍；当他行将熄灭时，他燃烧得最烈。他必然因赦免之恩使如此多死亡的工作在人身上被推翻，如此多先前属于他自己的定罪痕迹被抹去而哀伤叹息。他哀叹那罪人，如今是基督的仆人，注定要审判他和他的天使。\BibleRef{格前 6:3} 于是他观察、攻击、围困他，希望总能在某方面以肉体的贪欲击打他的眼目，或以世俗的诱惑缠住他的心思，或以属世权柄的恐惧颠覆他的信德，或以乖谬的传统使他离开正道：他从不缺少绊脚石和诱惑。因此，天主预见他的这些毒害，虽已藉圣洗的闩锁紧闭并封住了赦免之门，仍容许\Emph{它}微微开启。在门廊处，他安置了第二次补赎，为叩门者开门：但现在\Emph{一次}，因为第二次；但永不再有，因为上一次已属徒然。难道这仅\Emph{一次}还不够吗？你得到了你如今不配得的，因为你已失去了你所领受的。如果主的宽仁赐予你恢复所失之物的方法，你当感恩于这更新——不如说是加增的恩惠；因为\Emph{恢复}比\Emph{给予}更伟大，正如\Emph{失去}比从未\Emph{领受}更为悲惨。然而，若有人承担了第二次补赎的债务，他的灵魂不可立即被绝望所砍倒和破坏。\Emph{犯罪}再次固然可恶，但\Emph{悔改}再次却不可可恶：再次置身危险固然可恶，但再次获得解脱却不可。无人应感羞愧。屡次患病需屡次服药。你将因不拒绝主所赐之物而向主表达感恩。你已得罪，但尚能和好。你有一位可令其满意且他乐意者。\par

% structure: p0018 source=h2 target=section reason=relative-heading-rank; base=h2

\section{第8章 圣经中证明主乐意赦免的例证}

你若对此生疑，请细究\Quote{圣神向教会所说的话}之含义。祂归咎于厄弗所人\Quote{舍弃了起初的爱}\BibleRef{默 2:4}；责备提雅提辣人犯\Quote{奸淫}和\Quote{吃祭偶像之物}\BibleRef{默 2:20}；指控撒尔德人\Quote{行为不完善}\BibleRef{默 3:2}；谴责培尔加摩人教导邪僻之事\BibleRef{默 2:14}；斥责劳狄刻雅人倚靠自己的财富\BibleRef{默 3:17}；然而祂仍给他们所有人普遍的悔改劝告——固然是带着威吓；但若祂不宽恕悔改者，就不会向一个\Emph{不悔改}的人发出威吓。假若祂没有在别处彰显祂这种丰厚的仁慈，此事便仍属可疑。祂岂不是说：\Quote{那跌倒的必再起来，那\Emph{已转离的}必\Emph{转回}}？祂正是那一位\Quote{喜爱仁慈胜过祭献}者。诸天和其中的天使因一个人的悔改而喜乐\BibleRef{路 15:7}。嗨！你这罪人，振作起来吧！你看见那喜悦你归回的情景了。主的这些比喻对我们有何意义？一个妇女丢失了一块钱，便寻找它直到找到，并邀请女友共享喜悦，这不正是罪人归正的例子吗\BibleRef{路 15:8}？还有，牧羊人的一只小羊迷了路，但羊群并不比这一只更宝贵：那一只被切切寻找；那一只比所有更被渴慕；终于找到了，被牧羊人亲自扛在肩上归来，因为它在迷途中受了许多苦\BibleRef{路 15:3}。同样，那位最慈爱的父亲，我也不愿缄默不语：他召唤浪子回家，甘心乐意地接待了他贫困悔改的儿子，宰了最肥的牛犊，以盛宴点缀他的喜悦\BibleRef{路 15:11}。为何不呢？他找回了所失去的儿子；他因\Emph{得回}了\Emph{他}而更觉其宝贵。我们该理解这父亲是谁呢？无疑是天主：没有谁像祂那样真正是父亲；没有谁如此富于慈父之爱。那么，祂必会接纳你——祂自己的儿子——回家，即使你挥霍了从祂领受的一切，即使你赤裸着回来——仅仅因为你回来了；祂会为你的归回而喜悦，胜过为另一个儿子的持重而喜悦；但\Emph{唯有}当你真心悔改——若你将你的饥饿与你父亲家中\Quote{佣工}的丰裕相比——若你撇下猪群，那不清洁的畜群——若你再次寻求你的父亲，即使祂被冒犯，你说：\Quote{我犯了罪，不再配称为祢的儿子。}罪恶的告解减轻罪过，恰如掩饰加重罪过；因为告解出于（渴望）补赎，掩饰出于顽梗。\par

% structure: p0020 source=h2 target=section reason=relative-heading-rank; base=h2

\section{第九章 论伴随这第二次悔改的外在表现}

这第二次且唯一（剩下）的悔改，其行动范围越窄，考验就越艰巨；为使它不仅显于良心，且也落实于某种外在行为。这行为通常以希腊名称表达，即\OriginalTerm{悔改宣认}{\GreekText{ἐξομολόγησις}}，借此我们向主告明我们的罪过——并非似乎主不知道，而是因为告解得以确定补赎，忏悔由告解而生；藉忏悔天主得平息。因此，悔改宣认是一种使人谦卑、降伏的操练，要求一种能感动怜悯的举止。至于衣着和饮食，它命令（忏悔者）躺卧在苦衣和灰烬中，以衣服蒙羞，以痛苦使心灵谦卑，以严厉的对待来交换所犯的罪过；还规定只食用清淡的饮食——不是为了肚腹，而是为了灵魂；大抵上，以禁食喂养祈祷，叹息，哀哭，向主你的天主呼号；俯伏在长老们脚前，跪在天主所亲爱的人面前；嘱咐所有弟兄作使者，为他呈上恳求（在神前）。这一切悔改宣认（所做的），为的是增加悔改；因惧怕（所招致）的危险而尊崇天主；亲自定断罪人，代替天主的义怒，并藉暂时的禁欲（我不说挫败，而是）消除永罚。因此，它贬抑人，也高举人；它以污秽覆盖人，却使人更为洁净；它控告人，也辩护人；它定罪人，也赦免人。你对自己手下留情越少，（相信我）天主对你留情就越多。\par

% structure: p0022 source=h2 target=section reason=relative-heading-rank; base=h2

\section{第十章 论人对这第二次悔改和悔改宣认的退缩，以及这种退缩的不合理性}

然而，大多数人或者回避这工作，视之为公开暴露自己，或者日复一日拖延。我猜想，他们顾惜羞耻之心甚于救恩；就像那些身体隐秘部位患上疾病的人，躲避医生的诊察，因而死在自己的羞怯中。诚然，向被冒犯的主赔补，对羞耻之心来说是无法忍受的！要恢复已丧失的救恩！你的羞愧实在可敬；犯罪时脸皮厚，求饶时却羞答答！当羞耻的丧失使我获益时，我不给它留地步；当它本身在某儿子身上劝诫那人说：\Quote{不要顾惜我；我因你而丧亡，总比你因我而丧亡更好。}无论如何，当它（羞耻）的危险严重时，那是在嘲笑者面前成为戏谑对象的时候，这时一人以邻人的毁灭抬高自己，在跌倒者身上往上爬。但在弟兄和同仆中，有着共同的希望、恐惧、喜乐、忧伤、痛苦，因为来自共同的主和父的同一圣神，你为什么认为这些\Emph{弟兄}与你有什么不同？为什么要逃避你自己不幸的伙伴，好像他们会嘲笑你一样？身体不能因某个肢体的痛苦而高兴\BibleRef{格前 12:26}，它必须同心同意地忧伤，并努力寻求治疗。两个人在一处就是教会；而教会就是基督。因此，当你俯伏在弟兄们的膝前时，你是在触摸\Emph{基督}，你是在恳求\Emph{基督}。同样，当他们为你流泪时，是\Emph{基督}在受苦，是\Emph{基督}在向父祈求怜悯。儿子所求的总是容易得到。虚妄的羞耻之报酬何等伟大，它许诺我们隐藏过犯！也就是说，如果我们向人的知识隐藏一些事，难道就能同样向天主隐藏吗？难道人的审判与天主的知识是如此等同吗？暗中被定罪比公开被赦免更好吗？\Emph{但你说}：\Quote{这样来行\Emph{\OriginalTerm{公开忏悔}{exomologesis}}是可悲的。}是的，因为邪恶确实带来痛苦；但哪里悔改得以实行，痛苦就停止，因为它转化为有益的事。被切割、被烧灼、被某种药粉的刺激所折磨是可悲的：然而，那些以不快方式医治之物，因治疗之益处而自辩其可厌，并为了后续的好处而忍受当下的伤害。\par

% structure: p0024 source=h2 target=section reason=relative-heading-rank; base=h2

\section{第十一章 对同一主题的进一步批评}

此外，除了他们最在意的羞耻，难道\Emph{人们}不也惧怕身体上的不便吗？即他们必须不梳洗、衣着污秽、远离喜乐，在粗毛衣的粗糙和灰尘的可怕以及因禁食而凹陷的面容中度日？那么，我们穿着朱红和紫色来为我们的罪恳求合适吗？快拿发夹来整理头发，拿牙粉来抛光牙齿，拿一些钢或铜制的叉形工具来清洁指甲。无论有什么虚假的光彩、什么伪造的红润，\Emph{他都要}仔细地涂在嘴唇或脸颊上。此外，让他寻找花园或海边更温和温度的浴池；增加他的开销；仔细寻找最稀有的肥禽美味；精炼他的陈酒。当有人问他：\Quote{你把这一切都用在了谁身上？}让他说：\Quote{我得罪了天主，处于永远丧亡的危险中；所以我现在垂头丧气、消瘦、折磨自己，以求得与天主和好，因为我犯罪得罪了祂。}看哪，那些四处奔走谋求公职的人，为了他们的欲望而忍受灵魂和身体的烦恼，并不觉得可耻或厌烦；不仅仅是烦恼，还有各种侮辱。他们不装出什么低贱的衣着？不早早晚晚地拜访什么人家？——遇到任何显贵就鞠躬，不参加宴会，不参与娱乐，自愿放逐于自由和欢乐的福乐之外；所有这些，都是为了为期一年的短暂欢乐！当永恒处于危险时，\Emph{我们}难道还迟疑去忍受觊觎执政官或裁判官职位的人所忍受的吗？我们难道还迟迟不肯向被冒犯的主献上在食物和衣着上的自罚，而外邦人在没有得罪任何人的时候却对自己这样罚？圣经提到的正是这样的人：\Quote{祸哉，那些用长绳捆绑自己罪恶的人。}\par

% structure: p0026 source=h2 target=section reason=relative-heading-rank; base=h2

\section{第十二章 激励公开忏悔的最后考量}

你若畏缩于\Emph{\OriginalTerm{忏悔}{exomologesis}}，就当在心中思量那将被忏悔为你熄灭的地狱；先想象惩罚的严酷，好叫你不再迟疑采取这疗法。当那永火的小小通风口便能掀起如此烈焰，使邻近的城市或已不复存在，或每日担惊受怕时，我们又将那永恒之火的宝库看为何物？最傲然的山脉也因内在孕生的火在阵痛中裂开；而——这向我们证明审判的\Emph{永恒性}——尽管它们裂开，尽管它们被吞噬，却永无终结。谁不会将这些加于山脉的偶然惩罚视为那威胁不悔改者的审判的预表？谁不会同意，这些火花不过是从那不可估量的巨大火心中射出的些许飞镖和戏弄的矢石？因此，既然你知道在主圣洗的第一道壁垒之后，仍为你保留着\Emph{\OriginalTerm{忏悔}{exomologesis}}这对抗地狱的第二重援助，你为何要抛弃自己的救恩？为何迟迟不肯接近那深知能治愈你的疗法？甚至无理性的畜类也在急需时认出天主为它们指定的药草。被箭射穿的鹿，知道要拔出那铁镞及其难解难分的残留，必以白鲜自疗。燕子若弄瞎了幼雏，知道如何用自己的白屈菜使它们复明。罪人既知\Emph{\OriginalTerm{忏悔}{exomologesis}}是主为使他复原而设立的，岂可忽略那曾使巴比伦王恢复王权的事？他曾长久向主献上悔改，以七年的污秽执行他的\Emph{\OriginalTerm{忏悔}{exomologesis}}，指甲如鹰般疯长，乱发如狮般蓬乱。何等严酷的对待！众人所战栗的那一位，天主却重新接纳了他。反之，埃及皇帝——他曾追赶那曾经受折磨、长期抵制其主的天主之民，投入战争——在如此多的警告灾祸之后，却在分开的海中（那海只准\Quote{人民}通行）因波浪回流而灭亡：\BibleRef{出 14:15}因为他早已抛弃了悔改及其婢女\Emph{\OriginalTerm{忏悔}{exomologesis}}。\par

我何必要再多加赘述这人性救恩的两块木板（姑且这么说），更关心笔墨之事而非良心的责任？因为，我既是形形色色的罪人，生除悔改外别无他求，便不能轻易缄默于那连人类与人类过犯的根由和源头——亚当——也藉着\Emph{\OriginalTerm{忏悔}{exomologesis}}复归自己的乐园时不曾缄默的事。\par

\SourceRecord{Translated by S. Thelwall.；Ante-Nicene Fathers；Vol. 3.；Edited by Alexander Roberts, James Donaldson, and A. Cleveland Coxe.；Buffalo, NY: Christian Literature Publishing Co.,；1885.；<http://www.newadvent.org/fathers/0320.htm>.}
